
\subsubsection*{4.3.1}
\phantomsection
\label{sga3:VIIB:4.3.1}

With the notation of \SGARefLink{sga3:VIIB:4.3}{4.3}, let us determine the elements
\(\xi\) of \(U\) of augmentation \(1\) such that
\[
  \Delta_U(\xi)=\xi\otimes\xi:
\]
every element \(\xi\) of \(U\) can be written in the form
\[
  \xi
  =
  \sum_{n\in\mathcal P}
  \xi_n\frac{x^n}{n!},
\]
where \(\xi_n\in C\).
The condition \(\varepsilon_U(\xi)=1\) gives
\[
  \xi_0=1,
\]
where \(0\) denotes the element of \(\mathcal P\) all of whose
components vanish. The condition
\(\Delta_{U_p(L)}(\xi)=\xi\otimes\xi\) gives
\[
  \sum_{\substack{n\in\mathcal P\\m\leq n}}
  \xi_n
  \frac{x^m}{m!}
  \otimes
  \frac{x^{\,n-m}}{(n-m)!}
  =
  \sum_{q,r\in\mathcal P}
  \xi_q\xi_r
  \frac{x^q}{q!}
  \otimes
  \frac{x^r}{r!}.
\]
Equivalently,
\[
  \xi_{q+r}=\xi_q\xi_r
  \quad\text{if }q_i+r_i<p\text{ for every }i,
\]
and
\[
  \xi_q\xi_r=0
  \quad\text{otherwise}.
\]

Let \(\xi_i\) denote the coefficient \(\xi_n\) for which
\(n_i=1\) and \(n_j=0\) whenever \(j\ne i\). The preceding
conditions say precisely that
\[
  \xi_n=\prod_i\xi_i^{\,n_i}
  \quad(n\in\mathcal P),
  \qquad
  \xi_i^p=0
  \quad(i\in I).
\]
We obtain:

\medskip
\noindent\textbf{Proposition.}
\phantomsection
\label{sga3:VIIB:4.3.1:group-like-parametrization}
{\itshape Let \(p\) be a prime number, \(k\) a pseudocompact ring
of characteristic \(p\), \(\mathbf L\) a flat \(p\)-Lie algebra over
\(\mathbf O_k\), \(C\) a profinite \(k\)-algebra of finite length, and
\(\sqrt[p]{0_C}\) the ideal of \(C\) formed by the elements \(x\) such
that \(x^p=0\). There is a bijection
\[
  \mathbf L(C)\otimes_C\sqrt[p]{0_C}
  \longrightarrow
  \mathscr G_p(\mathbf L)(C)
\]
``which is functorial in \(C\).''}\footnote{Translator's note: the Gille--Polo re-edition (2011),
Expos\'e~VII\textsubscript{B}, assumes \(k\) local, so that the projective pseudocompact \(k\)-module dual to
\(\mathbf L\) is topologically free (0.2.1), as the choice of coherent bases in the proof below (1.2.3)
requires. The statement holds without this hypothesis: \(k\) is the product of its local components \(k_m\)
(0.2.1); every \(C\) in \(\operatorname{Alf}/k\) is the product of the finitely many nonzero
\(C_m=k_m\otimes_kC\), and \(\mathbf L(C)\otimes_C\sqrt[p]{0_C}\), \(U_p(\mathbf L(C))\) and
\(\mathscr G_p(\mathbf L)(C)\) are the products of the corresponding objects for the \(C_m\); the bijections
given by the proof over each \(k_m\) are compatible with the morphisms of \(\operatorname{Alf}/k\), which
respect these decompositions, and their product is the required bijection. Theorem 4.4 uses the statement for
an arbitrary \(k\).}

\textit{Proof.}
By \SGARefLink{sga3:VIIB:1.2.3}{1.2.3}, one may indeed choose a basis
\[
  ({}^C x_i)_{i\in I}
\]
of \(\mathbf L(C)\) over \(C\) in such a way that, for every
morphism \(\phi:C\to D\) of \(\operatorname{Alf}/k\)
\textup{(\SGARefLink{sga3:VIIB:0.4.2}{0.4.2})}, the map \(\mathbf L(\phi)\) sends \({}^C x_i\) to
\({}^D x_i\). By what precedes, the map
\[
  \sum_i({}^C x_i)\otimes\xi_i
  \longmapsto
  \sum_{n\in\mathcal P}
  \left(\prod_i\xi_i^{\,n_i}\right)
  \frac{({}^C x)^n}{n!}
\]
is a functorial bijection from
\(\mathbf L(C)\otimes_C\sqrt[p]{0_C}\) onto
\(\mathscr G_p(\mathbf L)(C)\).

